\documentclass[11pt]{article}
\usepackage[a4paper,margin=1in]{geometry}
\usepackage[T1]{fontenc}
\usepackage{lmodern}
\usepackage[hidelinks]{hyperref}
\usepackage{parskip}
\pagestyle{empty}
\begin{document}
\begin{flushright}
Quang-Vinh Dang\\
British University Vietnam\\
Ecopark Township, Hung Yen, Vietnam\\
\href{mailto:vinh.dq4@buv.edu.vn}{vinh.dq4@buv.edu.vn}\\
ORCID: \href{https://orcid.org/0000-0002-3877-8024}{0000-0002-3877-8024}\\[1em]
30 September 2026
\end{flushright}

The Editor-in-Chief\\
\emph{Journal of Digital Economy}

Dear Editor,

I am submitting the manuscript ``Searching for the digital dividend: Internet diffusion and productivity growth in 149 economies, 1996--2025'' for consideration as a research article in the \emph{Journal of Digital Economy}.

The paper asks whether the spread of internet use has been followed by faster labour-productivity growth, and whether any such association differs across economies. It uses a panel of 149 economies over 1996--2025 built from the World Development Indicators, the UNDP Human Development Report data and ITU series, with Penn World Table 10.0 and 11.0 as alternative sources for productivity.

The main findings are as follows:
\begin{itemize}\setlength\itemsep{0pt}
\item Within countries, years with a higher internet-user share were not years of faster productivity growth. A linear fixed-effects estimate of $-0.28$ percentage points per ten-point higher share is robust to country-specific trends and convergence controls. A cross-fitted machine-learning estimate is less precise, and its interval includes zero under two-way clustering. The association is concentrated in 1996--2011.
\item The association does not vary robustly with schooling, initial productivity or time since internet take-off, although the data cannot rule out moderate gradients.
\item A gradient along initial productivity appears in both Penn World Table releases but is much weaker in World Bank data for the same country-years. About half of the difference comes from the productivity-growth series.
\end{itemize}

To estimate the heterogeneity, the paper develops Panel-DOSE, a debiased machine-learning spline estimator for panels with two-way fixed effects and dependence within countries. The paper derives its main properties and reports simulations, including a plasmode design built on the empirical panel, that show where the estimator helps and where simpler methods do as well.

I believe the paper fits the journal's scope on the economic consequences of digital technologies. It provides cross-country evidence on the productivity paradox of the internet era, and it shows how the choice of data source can change conclusions about who benefits from digital adoption. The results are observational and are presented as conditional associations, not as causal returns to connectivity investment.

The manuscript is original. It has not been published and is not under consideration for publication elsewhere. The author declares no competing interests, and the research received no specific funding. The study uses only publicly available country-level data and did not require ethics approval. Generative AI tools were used to assist with code and with the writing, as stated in the declaration on the title page; the author reviewed and edited all content and takes full responsibility for it. An anonymised replication package with the data, code and scripts that reproduce every result is provided with the submission.

The submission consists of the anonymised manuscript, the anonymised supplementary material, a separate title page, the highlights, the ethics statement and the replication package.

Thank you for considering this manuscript. I look forward to hearing from you.

Yours sincerely,

Quang-Vinh Dang\\
British University Vietnam
\end{document}
